% compendium/execution-trading-costs/transaction-cost-analysis.tex
% Entries migrated on 2026-10-15 from the legacy transaction-cost-analysis file (slice 13), owned by the
% section overview execution-trading-costs; dispositions and source checks in notes/migration.md,
% the blocks verbatim in notes/artefacts/market-making-lob/derivations.tex.

\section{Transaction costs and transaction cost analysis}

\begin{qfdefinition}[Transaction costs: explicit, implicit and missed-trade costs, and an accounting split of shortfall]\label{def:transaction-costs}
\qffield{Statement} Transaction costs include all costs associated with trading \citep[chap.~``Liquidity and Transaction Cost Measurement'', sec.~21.1, pp.~421--422]{alpha-harris}:
\begin{itemize}
\item \emph{explicit costs}, which a cost accountant would identify: commissions, exchange fees, taxes, and the resources devoted to trading, such as running a buy-side desk;
\item \emph{implicit costs}, which arise because traders affect prices: the bid--ask spread paid by buying at the ask and selling at the bid, and the price impact of large orders;
\item \emph{missed trade opportunity costs}, when orders fail to fill, or fill late. Harris's example: a cotton buyer who bid 64.95 cents when the price was 65 and saw it rise to 68 lost 2.75 cents a pound against an aggressive fill at 65.25, \$137{,}500 on 100 contracts (checked).
\end{itemize}
Explicit costs are easy to measure. Implicit and missed-trade costs need a benchmark: what prices would have been without the trade, or what the unfilled order would have cost to complete (p.~422). Kissell lists ten components, commissions, fees, taxes, rebates, spreads, delay cost, price appreciation, market impact, timing risk and opportunity cost, and classifies them as investment-related, trading-related and opportunity costs, fixed or variable, visible or hidden \citep[chap.~3, pp.~89--95, Tables~3.1--3.2]{s7-kissellscience}.

\qffield{Notation} Interpretation: with the signed quantities of \hyperref[def:implementation-shortfall]{implementation shortfall} ($S$ intended, $s_j$ filled at $p_j$, decision price $P_d$, arrival midpoint $P_0$, end price $P_n$) and $m_j$ the quotation midpoint just before fill $j$, shortfall splits additively as
\[ \mathrm{IS}=\underbrace{S(P_0-P_d)}_{\text{delay}}+\underbrace{\textstyle\sum_js_j(p_j-m_j)}_{\text{half-spread paid}}+\underbrace{\textstyle\sum_js_j(m_j-P_0)}_{\text{midpoint drift}}+\underbrace{\bigl(S-\textstyle\sum_js_j\bigr)(P_n-P_0)}_{\text{opportunity}}+\text{fees}. \]
The second term is the filled size times Harris's liquidity premium, half the effective spread at each fill (\hyperref[def:spread-measures]{spread measures}); the middle two terms are Wagner's trading-related part. Checked with exact arithmetic: a buy of 10{,}000 shares (decision 50.00, arrival 50.10, fills of 3{,}000 at 50.20 and 50.30 against pre-fill midpoints 50.15 and 50.25, rest cancelled, $P_n=50.60$, fees 60) gives $1{,}000+300+600+2{,}000+60=3{,}960$, the paper-minus-real total; the mirror-image sell gives the same.

\qffield{Conditions} The split is an accounting identity, not a causal attribution. The midpoint drift mixes the order's own impact with market movement and other traders' impact. Kissell defines market impact as the difference between the price path with the order and without it, which cannot be observed together, so impact cannot be measured exactly \citep[chap.~3, pp.~91--92]{s7-kissellscience}; Harris likewise needs a no-trade benchmark (p.~422). The half-spread term measures each fill against its own contemporaneous midpoint, so it misses the cumulative impact of earlier fills, which falls into the drift (Harris's split-order example, pp.~428--429). Kissell's timing risk is the uncertainty of cost, not a realised component (p.~92), and his price appreciation is the price trend that would have occurred anyway (p.~91).

\qffield{Practice} Interpretation: an index or market adjustment (\hyperref[meth:transaction-cost-analysis]{transaction cost analysis}) removes the part of the drift predicted by the market's move; what remains is still a residual, which includes impact, stock-specific news and model error, and is impact only under the model used to remove the rest. The overview's checkpoint asks for a split into spread, impact, timing and fees: on the sources checked, spread, delay (timing before arrival), opportunity and fees are measurable as above, but ``impact'' is available only as a model estimate or a market-adjusted residual, not as an observed component.

\qffield{Pitfalls} Count each cost once: the spread is inside fill prices, so adding a quoted or effective spread to an arrival-price shortfall counts it twice; fees enter once. The legacy taxonomy listed ``adverse selection'' and ``delay cost'' among implicit costs ``relative to a benchmark''; in Harris adverse selection is a cause of the spread, and in Kissell delay is an investment-related cost before the order reaches the market. ``Borrow charges'' among explicit costs is not in the passages checked. Kissell describes the spread as the round-trip cost for small orders, not blocks (p.~90).

\qffield{See also} \hyperref[def:implementation-shortfall]{implementation shortfall}; \hyperref[def:execution-benchmark]{execution benchmarks}; \hyperref[def:spread-measures]{spread measures}, an entry of the microstructure overview; \hyperref[def:price-impact]{price impact}, an entry of the price-impact pathway; \hyperref[meth:transaction-cost-analysis]{transaction cost analysis}.

\qffield{Further reading} The legacy blocks it replaces, verbatim: \nolinkurl{notes/artefacts/market-making-lob/derivations.tex}.

\qffield{Index} transaction costs; explicit costs; implicit costs; missed trade opportunity cost; commissions; fees; taxes; rebates; spread; delay cost; price appreciation; market impact; timing risk; opportunity cost; investment-related cost; trading-related cost; midpoint drift; liquidity premium; cost decomposition

\qffield{Status} Verification: partly checked 2026-10-15 (against the text layers of local copies of \texttt{alpha-harris}, 2002 Oxford printing, pp.~420--423 and 428--429, and \texttt{s7-kissellscience}, Academic Press printing with copyright 2014, chap.~3, pp.~87--96; the split and examples checked with exact arithmetic). Stability: draft. Currency: timeless (accounting).
\end{qfdefinition}

\begin{qfmethod}[Transaction cost analysis: measuring costs after the trade, predicting them before, and comparing brokers and algorithms]\label{meth:transaction-cost-analysis}
\qffield{Statement} Traders measure transaction costs to evaluate how well they and their brokers implemented their strategies, whether they traded too aggressively or not enough, and which brokers should get future orders; they predict costs to judge whether strategies are worth implementing. Exchanges, brokers (notably to show best execution when paid for order flow), sponsors (ERISA commission duties) and regulators measure them too \citep[chap.~``Liquidity and Transaction Cost Measurement'', pp.~420--421]{alpha-harris}. Lehalle and Laruelle separate post-trade analysis, usually of one day, which explains what the market context did, from TCA over months of trades, which uses averaging to cancel most market effects and find what part of performance comes from the algorithm's design; its outcome is to match an algorithm or provider to a market context or trading style, so that flow is routed accordingly, using breakdowns by volatility, jumps, liquidity and news, and comparing only stocks of similar liquidity (bid--ask spread, tick size in basis points, quantity on the book, turnover, trade size, free float, trading rate) \citep[sec.~3.1.3, pp.~202--203]{alpha-laruelle}.

\qffield{Notation} Kissell's measures \citep[chap.~3, pp.~110--113]{s7-kissellscience}: index cost $=(\text{index VWAP}-\text{index arrival})/\text{index arrival}\times10^4$ over the trading horizon (eq.~(3.21)); index-adjusted cost $=$ arrival cost $-\hat\beta\times$ index cost (3.22); and $Z=(\text{pre-trade estimate}-\text{arrival cost})/\text{pre-trade timing risk}$ (3.23), positive when the trade beat its estimate. A cost \emph{measure} is ex post and single-valued; a cost \emph{forecast} is ex ante, a mean with a standard deviation, the timing risk (pp.~95--96).

\qffield{Conditions} Prediction: explicit information (contracted fees and displayed liquidity, an upper bound on impact for orders smaller than displayed size) and implicit information from past costs, modelled by regression on order variables (size, price placement), contemporaneous market variables (spread, displayed size, recent volume, recent returns, money flow) and general ones (average volume, volatility, market capitalisation); past costs predict future ones only if liquidity conditions are unchanged \citep[chap.~``Liquidity and Transaction Cost Measurement'', sec.~21.7, pp.~438--439]{alpha-harris}. Comparison: estimates from single trades are noisy, averaging removes noise but not bias or gaming, and comparing brokers is meaningful only for similar trading problems (p.~432). Kissell's $Z$ should be near mean zero and variance one if pre-trade estimates are accurate; the test should be run by order size and strategy, on shares actually traded rather than the full order, on enough orders, and with the investor's own pre-trade model rather than the broker's (pp.~112--113).

\qffield{Practice} Kissell separates cost from performance: comparing the execution price with the day's VWAP is a performance proxy, with the close a tracking-error proxy, while comparison with the open or the arrival price is a cost (p.~96). Harris advises that strategists and traders share why they trade, and that incentives reward joint performance rather than blame (pp.~439--440).

\qffield{Pitfalls} Interpretation: the population, aggregation weights and benchmark must be stated: averaging by order and by value answer different questions, and both depend on which orders reached each broker. Lehalle and Laruelle's averaging cancels market effects only if they are unrelated to which algorithm was chosen; if traders send urgent or difficult orders to one provider, its average cost is worse whatever its quality, so a routing table built on unadjusted averages confounds selection with skill, and outperformance of a benchmark does not show causal superiority. The legacy failure modes (wrong exposure, wrong clock, own-impact contamination, gaming, selection) are consistent with Harris's biases (sec.~21.4) and Kissell's conditions but are a synthesis, not a list in either. Kissell's index cost is ``not adjusted for the side of the trade'' (p.~112); with a side-signed arrival cost this is wrong for sells: a sale at 50~basis points below arrival while the index fell 50~basis points and $\hat\beta=1$ shows 100~basis points of adjusted cost instead of 0, so the index term needs the side factor (Interpretation). His approximation ``Index VWAP $=\tfrac12R_m$'' is an approximation to the index cost, not the index VWAP. The legacy prediction variables added venue variables (hidden fill rates, fee tiers, queue loss, routing latency) and limit price and urgency, which are not in the passages checked.

\qffield{See also} \hyperref[def:transaction-costs]{transaction costs}; \hyperref[def:execution-benchmark]{execution benchmarks}; \hyperref[meth:execution-decision-support]{pre-trade, monitoring and post-trade tools}; \hyperref[meth:aggressiveness-tradeoff]{aggressiveness trade-off}; \hyperref[def:best-execution]{best execution}; the market adjustment parallels \hyperref[meth:relative-performance-evaluation]{relative performance evaluation} of managers.

\qffield{Further reading} The legacy blocks it replaces, verbatim: \nolinkurl{notes/artefacts/market-making-lob/derivations.tex}.

\qffield{Index} transaction cost analysis; TCA; post-trade analysis; cost measurement; cost prediction; pre-trade cost model; index-adjusted cost; market-adjusted cost; z-score; timing risk; broker evaluation; algorithm selection; routing table; selection bias; liquidity breakdowns; cost versus performance

\qffield{Status} Verification: partly checked 2026-10-15 (against the text layers of local copies of \texttt{alpha-harris}, 2002 Oxford printing, pp.~420--421, 432 and 438--440; \texttt{alpha-laruelle}, 2018 World Scientific printing, pp.~202--203; and \texttt{s7-kissellscience}, Academic Press printing with copyright 2014, chap.~3, pp.~95--96 and 110--113, p.~112 by page image). Stability: draft. Currency: practitioner methods as of 2002, 2014 and 2018.
\end{qfmethod}

\begin{qfmethod}[Trading aggressiveness: balancing the transaction cost of trading harder against the missed-trade cost of trading softer]\label{meth:aggressiveness-tradeoff}
\qffield{Statement} If the additional transaction costs of trading more aggressively are less than the missed trade opportunity costs it would save, traders earn more by trading more aggressively; if the transaction costs saved by trading less aggressively exceed the additional opportunity costs, they earn more by trading less aggressively. Most traders do not estimate these marginal costs but assume they equal average costs, trading harder or softer as their per-unit opportunity costs exceed or fall short of their per-unit transaction costs; Harris advises estimating marginal transaction costs from the costs of the last trades that fill orders \citep[chap.~``Liquidity and Transaction Cost Measurement'', sec.~21.6, pp.~437--438]{alpha-harris}.

\qffield{Conditions} Interpretation: the objective is expected total cost, transaction cost plus missed-trade cost, of an order whose size is given; the choice variable is aggressiveness (limit placement, urgency, participation); the rule is the first-order condition of that minimisation and assumes both costs change smoothly with aggressiveness and that the missed-trade cost is measured realistically. Harris states it as advice, with no model or proof. Opportunity costs count only for orders that could have been filled, and depend on the horizon at which they are measured (pp.~435--436).

\qffield{Practice} Well-informed traders, especially on short-term information, should accept substantial transaction costs rather than fail to trade, and strategists should give traders price limits that reflect the value of their information (p.~439). Interpretation: a mean--variance schedule (\hyperref[mod:almgren-chriss-execution]{Almgren--Chriss}) is one formal version of the same trade-off, with price risk in place of missed-trade cost.

\qffield{Pitfalls} The legacy note ended ``average cost is the wrong control variable for the last slice''; Harris's point is that marginal, not average, costs should drive the decision, and that the last fills estimate them. Interpretation: comparing average costs gives the right direction only if marginal and average costs move together.

\qffield{See also} \hyperref[def:transaction-costs]{transaction costs}; \hyperref[def:implementation-shortfall]{implementation shortfall} (the opportunity term); \hyperref[mod:almgren-chriss-execution]{mean--variance execution}; \hyperref[def:order-exposure-decisions]{order exposure decisions}.

\qffield{Further reading} The legacy block it replaces, verbatim: \nolinkurl{notes/artefacts/market-making-lob/derivations.tex}.

\qffield{Index} trading aggressiveness; marginal transaction cost; marginal opportunity cost; missed trade opportunity cost; average cost; order submission; urgency

\qffield{Status} Verification: partly checked 2026-10-15 (against the text layer of a local copy of \texttt{alpha-harris}, 2002 Oxford printing, pp.~435--439). Stability: draft. Currency: timeless (decision rule).
\end{qfmethod}
